\section{Evaluation}\label{sec:experiments}

Our evaluation aims to address three questions: (i) does \textsc{VOILA} consistently improve accuracy-cost trade-offs versus fixed-fidelity baselines? (ii) do gains persist across model scales and datasets? and (iii) which components are performance-critical?
\fablenote{Question (i) sets the bar too low --- fixed-fidelity is the weakest comparison class. Rewrite as ``versus fixed-fidelity policies and adaptive routing baselines (cascades, LLM/BERT routers)''; you already beat those baselines, so claim the harder fight up front. Also add a missing experimental detail AAAI reviewers always check: how many evaluation queries per dataset, how many runs/seeds, and whether differences are statistically significant (bootstrap CIs over queries cost nothing to add and immunize every comparison in \S\ref{sec:results}).}

\textbf{Models.} We evaluate six VLMs: Pixtral-12B~\cite{pixtral2024}, LLaMA-4-Maverick (17B$\times$12E)~\cite{abdullah2025evolutionmetasllamamodels}, LLaVA-1.5-7B~\cite{zhou2024tinyllava}, Qwen2-VL-72B~\cite{Qwen2VL}, Qwen2.5-VL-7B~\cite{Qwen2VL}, and Qwen3-VL-235B~\cite{bai2023qwen}, spanning 7B--235B parameters across architectures from 2023--2025. All models are used zero-shot.

\textbf{Datasets.}
\textbf{VQA-v2}~\cite{goyal2017vqav2} and \textbf{GQA}~\cite{hudson2019gqa} are standard in-domain VQA benchmarks.
\textbf{TextVQA}~\cite{singh2019textvqa} emphasizes OCR-dependent reasoning.
\textbf{FloodNet}~\cite{rahnemoonfar2021floodnet} is a bandwidth-constrained cyber-physical setting (aerial disaster imagery).
\textbf{LoCoMo}~\cite{wang2024locomo}: we generate VQA-style QA pairs from conversational memory items using GPT-4o (details in Appendix~\ref{sec:datasets}).

\textbf{Input Fidelities.}
$\mathcal{F}=\{\texttt{caption},\ \texttt{resize}_{32\times32},\ \texttt{jpeg\_q1},\ \texttt{jpeg\_q10},\ \texttt{full}\}$, ordered by acquisition cost. Captions are generated offline by the target VLM from full-resolution images; caption retrieval cost (7.6) reflects serving a pre-computed string from cache, not one-time generation. Full-resolution cost is normalized to 100 (details in Appendix~\ref{sec:cost_model}).

\paragraph{Cost Model.}
Each fidelity incurs acquisition cost $c(f) = \alpha B(f) + \beta L(f) + \gamma S(f)$, where $B(f)$, $L(f)$, $S(f)$ are empirically measured bandwidth, retrieval latency, and storage-access ratios. Costs reflect input acquisition overhead only; \textsc{VOILA} depends only on relative cost differences.

\textbf{Baselines.}
We compare \textsc{VOILA} against:
\begin{itemize}[nosep]
    \item \textbf{Fixed-fidelity} ($f \in \mathcal{F}$): always retrieve at a single level. Establishes the non-adaptive Pareto frontier.
    \item \textbf{Log-prob Cascade}~\cite{chen2023frugalgpt}: execute VLM at lowest fidelity; escalate if log-probability falls below threshold $\theta$ (tuned per model on VQA-v2 validation). Incurs retrieval and inference costs at every intermediate fidelity.
    \item \textbf{GPT-4o Router} (zero-shot): GPT-4o classifies each question's visual need as low/medium/high (mapped to caption / JPEG Q10 / full). Labels collected once per dataset, shared across all models.
    \item \textbf{BERT Router}: supervised SentenceTransformer (all-MiniLM-L6-v2) predicting the cheapest correct fidelity label directly --- no VOI objective.
    \item \textbf{Oracle}: cheapest fidelity that actually produces a correct answer (hindsight upper bound).
    \item \textbf{\textsc{VOILA}-CDF}: \textsc{VOILA} with VOI replaced by a CDF threshold rule. Ablates cost-aware routing while holding predictor and calibration fixed.
\end{itemize}

\textbf{Evaluation Metrics.}
We report task accuracy (\%) and normalized average retrieval cost per query (full-resolution = 100). For interpretability, we also report the fraction of queries routed to each fidelity level and the accuracy--cost Pareto frontier.
For calibration quality we additionally report Brier score and Expected Calibration Error (ECE) in the ablation studies (\S\ref{sec:results}).
\fablenote{Two-number (Acc, Cost) comparisons make it impossible to rank methods when neither dominates --- exactly the situation in Table~\ref{tab:master_comparison} where the GPT-4o Router has higher accuracy but higher cost than \textsc{VOILA} on VQA-v2. Add a single scalar metric that resolves this: either (a) area under the accuracy--cost curve as $\lambda$ sweeps, or (b) accuracy at matched cost / cost at matched accuracy. Without one, a reviewer can legitimately read Table~1 as ``the LLM router is more accurate.''}
